\documentclass[10pt,a4paper]{amsart}
\usepackage[margin=19mm]{geometry}
\usepackage{amsmath,amssymb,mathtools}
\usepackage[scr=rsfs,cal=euler]{mathalfa}
\usepackage{xfrac}
\usepackage[unicode,hidelinks,bookmarks=false]{hyperref}
\newcommand{\ie}{i.e.\ }
\newcommand{\cf}{cf.\ }
\newcommand{\resp}{respectively\ }
\newcommand{\Subsection}{\subsection}
\providecommand{\nobreakdash}{\nobreak\hbox{-}}
\setlength{\emergencystretch}{2em}
\multlinegap0pt
\usepackage{bm}
\usepackage{bbm}
\usepackage{dsfont}
\usepackage{xspace}
\usepackage{cancel}
\usepackage{mathtools}
\usepackage{amsthm}
\makeatletter
\@ifundefined{thm}{\newtheorem{thm}{Thm}}{}
\@ifundefined{lem}{\newtheorem{lem}[thm]{Lemma}}{}
\@ifundefined{prop}{\newtheorem{prop}[thm]{Proposition}}{}
\@ifundefined{rem}{\newtheorem{rem}[thm]{Remark}}{}
\makeatother
\title[Multifractality and intermittency in the limit evolution of ]{Multifractality and intermittency in the limit evolution of polygonal vortex filaments}
\author{Valeria Banica, Daniel Eceizabarrena, Andrea R. Nahmod, Luis Vega}
\date{Source: arXiv:2309.08114v4 (CC BY 4.0)}
\begin{document}
\maketitle

\noindent\emph{Source and licence.} Multifractality and intermittency in the limit evolution of polygonal vortex filaments. Version arXiv:2309.08114v4 (CC BY 4.0), distributed under the Creative Commons Attribution 4.0 International licence, as recorded in the arXiv licence metadata for this exact version and shown on the abstract page https://arxiv.org/abs/2309.08114v4. Full rights notice: licenses/arxiv-2309.08114-CC-BY-4.0.txt.\par\smallskip

\noindent\textit{Editorial scope.} Counted unit: the Prop whose source label is \texttt{thm:UpperBound34} in the article (source0.tex lines 3062--3064, source label \texttt{thm:UpperBound34}), delivered together with the article's own complete original proof of it (source0.tex lines 3065--3337, one \texttt{proof} environment of 273 lines). No mathematics is written, repaired or re-derived: statement and proof are the article's own text line for line. Required context retained uncounted: eq:Asymptotic\_xq\_Equal\_0 (lines 2981--2988); eq:Asymptotic\_xq\_Not\_0 (lines 2990--2997). Every definition, display and label of those lines is retained unchanged. Omitted from this package and not redistributed: 1--2980, 2989--2989, 2998--3061, 3338--3529. Nothing inside a retained excerpt is omitted or altered: every retained body is delivered exactly as the article's lines are written, so no omission receipt of any kind accompanies this package. Citations: the retained excerpts carry no citation command at all, so no bibliography entry of the article is transcribed and no reference is redistributed. Reference adaptation: none was needed and none was made; every reference inside the delivered unit resolves inside it, so no unresolved reference or citation is produced. Printed slips kept literal: no printed word, symbol, hypothesis or number of the counted statement or of its proof was altered. Layout and class: the article's own class is \texttt{amsart} with packages including amsmath, amssymb, amsthm, amsfonts, bbm, graphicx, pstricks, hyperref, url, tasks, enumerate, caption, subcaption, mathtools; the wrapper is the lane's pinned \texttt{amsart} editorial wrapper at 10pt on A4, which restates the theorem-like environments the retained text needs and adds the article's own macro definitions that the retained text needs (none), with extra wrapper packages (bm, bbm, dsfont, xspace, cancel, mathtools). The delivered numbering is the unit's own. Assets: no retained span contains a figure environment, a graphics-inclusion command, a PDF-page inclusion, a table or a TikZ picture; no image file is copied into this package, the package declares no asset dependency and no image file is redistributed with it. Licence: version arXiv:2309.08114v4 of this article is distributed under the Creative Commons Attribution 4.0 International licence, as recorded in the arXiv licence metadata for this exact version and shown on the abstract page https://arxiv.org/abs/2309.08114v4. A full rights notice is delivered at licenses/arxiv-2309.08114-CC-BY-4.0.txt. Characters: every retained excerpt is pure ASCII, so the delivered engine, XeLaTeX with its default Latin Modern text font, reports no missing character.
	\begin{equation}\label{eq:Asymptotic_xq_Equal_0}
\begin{split}
&R_{x_0}\Big( \frac{p}{q} + h \Big) - R_{x_0}\Big( \frac{p}{q} \Big) + 2\pi i h   \\
& \qquad   \quad = 2\pi (- 1 \pm i) \frac{\sqrt{|h|}}{\sqrt{q}} \frac{ G(p,m_q,q)}{\sqrt{q}} 
+  2 (1 \pm i) \, q^{3/2} |h|^{3/2} \,\sum_{m \neq 0} \, \frac{ G(p,m_q + m,q)}{\sqrt{q}} \,  \frac{e^{- 2\pi i \frac{m^2}{4 q^2 h}}}{ m^2}
 + O\left( q^{7/2} h^{5/2}  \right),
 \end{split}
\end{equation}

	\begin{equation}\label{eq:Asymptotic_xq_Not_0}
\begin{split}
& R_{x_0}\Big( \frac{p}{q} + h \Big) - R_{x_0}\Big( \frac{p}{q} \Big) + 2\pi i h  \\
& \qquad \qquad 
= 2 (1 \pm i)  \, q^{3/2} |h|^{3/2} \,\sum_{m \in \mathbb Z} \, \frac{ G(p,m_q + m,q)}{\sqrt{q}} \,  
\frac{e^{- 2\pi i \frac{(m - qx_q)^2}{4 q^2 h}}}{ (m - qx_q)^2} \,  + O\left( q^{7/2} h^{5/2} \sum_{m \in \mathbb Z} \frac{1 }{(m - qx_q)^4}  \right).
\end{split}
\end{equation}

\begin{prop}\label{thm:UpperBound34}
Let $x_0 \in \mathbb Q$ and $t \not\in \mathbb Q$. Then, $\alpha_{x_0}(t) \leq  3/4$. 
\end{prop}

\begin{proof}
Set $x_0 = P/Q$ with $P,Q \in \mathbb N$ and $(P,Q) = 1$. 
Let $t \not\in \mathbb Q$ and let $p_n/q_n$ be its approximations by continued fractions. 
It is well-known\footnote{Because two consecutive denominators $q_n$ and $q_{n+1}$ are never both even.} that there is a subsequence of odd denominators $q_{n_k}$. 
Renaming that subsequence back to $q_n$, we may assume that all $q_n$ are odd. 
Consequently, $|G(p_n, m, q_n)| = \sqrt{q}$ for all $m, n \in \mathbb N$. 
As usual, let 
\begin{equation}
h_n =  t - \frac{p_n}{q_n}, 
\qquad|h_n| < \frac{1}{q_n^2}, 
\qquad x_{q_n} = \min_{m \in \mathbb Z} \Big| \frac{P}{Q} - \frac{m}{q_n} \Big| = \Big| \frac{P}{Q} - \frac{m_{q_n}}{q_n} \Big|,
\end{equation}
and we immediately deduce that either $x_{q_n} = 0$ or $1/Q \leq q_n x_{q_n} \leq 1/2$. 
We separate cases:
\begin{itemize}
	\item[\textbf{Case 1}] We have $x_{q_n} = 0$ for infinitely many $n \in \mathbb N$. 
	Rename that subsequence and rewrite \eqref{eq:Asymptotic_xq_Equal_0} as 
	\begin{equation}\label{eq:Asymptotic_Case_1}
\Big| R_{x_0} \Big( \frac{p_n}{q_n} + h_n \Big)  - R_{x_0}\Big( \frac{p_n}{q_n} \Big) + 2\pi i h_n   \Big| 
 = 2\pi \sqrt{2} \frac{\sqrt{|h_n|}}{\sqrt{q_n}}  + O\left( q_n^{3/2} h_n^{3/2}  \right) 
 \simeq \frac{\sqrt{|h_n|}}{\sqrt{q_n}} \left(  1 + O\big( q_n^2 h_n \big) \right).
\end{equation}
Let $\delta > 0$ which we determine later. Separate cases again:
\begin{itemize}
	\item[\textbf{Case 1.1.}] Suppose that $\left|  1 + O\big( q_n^2 h_n \big) \right| \geq \delta$ for infinitely many $n \in \mathbb N$. 
	Then, 
	\begin{equation}
	\left| R_{x_0} (t) - R_{x_0}( t - h_n ) + 2\pi i h_n   \right| 
	\geq \delta \frac{\sqrt{|h_n|}}{\sqrt{q_n}}  
	\geq \delta \, |h_n|^{3/4}, 
	\end{equation}
	because $q_n^2 |h_n| \leq 1$. 
	Hence $| R_{x_0} (t) - R_{x_0}( t - h_n )| \geq (\delta/2) \, |h_n|^{3/4}$ for infinitely many $n \in \mathbb N$,
	and consequently $\alpha_{x_0}(t) \leq 3/4$. 
	
	\item[\textbf{Case 1.2.}] We have $\left|  1 + O\big( q_n^2 h_n \big) \right| < \delta$ for all large enough $n$.
	In that case, we evaluate \eqref{eq:Asymptotic_Case_1} at a point closer to $p_n/q_n$.
	Let $\epsilon > 0$ and write \eqref{eq:Asymptotic_xq_Equal_0} for $\epsilon h_n$, so that instead of  \eqref{eq:Asymptotic_Case_1} we get
	\begin{equation}
\Big|   R_{x_0} \Big( \frac{p_n}{q_n} + \epsilon h_n \Big)  - R_{x_0}\Big( \frac{p_n}{q_n} \Big) + 2\pi i \epsilon h_n \Big|   
 \simeq \sqrt{\epsilon} \, \frac{\sqrt{|h_n|}}{\sqrt{q_n}} \left(  1 + \epsilon O\big( q_n^2 h_n \big) \right).
\end{equation}
Since $q_n^2 |h_n| < 1$ and the constant underlying the big-$O$ is universal, say $C$, choose $\epsilon \leq 1/(2C)$, in such a way that 
\begin{equation}
\Big| R_{x_0} \Big( \frac{p_n}{q_n} + \epsilon h_n \Big)  - R_{x_0}\Big( \frac{p_n}{q_n} \Big) + 2\pi i \epsilon h_n  \Big|   
\gtrsim  \frac{\sqrt{\epsilon}}{2} \, \frac{\sqrt{|h_n|}}{\sqrt{q_n}}.
\end{equation}
From this and \eqref{eq:Asymptotic_Case_1}, we write
\begin{equation}
\begin{split}
\frac{\sqrt{\epsilon}}{2} \, \frac{\sqrt{|h_n|}}{\sqrt{q_n}} 
& \lesssim \Big| R_{x_0} \Big( \frac{p_n}{q_n} + \epsilon h_n \Big)  - R_{x_0}\Big( \frac{p_n}{q_n} \Big) \Big| + 2\pi \epsilon |h_n|  \\
& \leq \Big| R_{x_0} \Big( \frac{p_n}{q_n} + \epsilon h_n \Big)  - R_{x_0}( t ) \Big| + \Big| R_{x_0} ( t )  - R_{x_0}\Big( \frac{p_n}{q_n} \Big) \Big| + 2\pi \epsilon |h_n| \\
& \lesssim \Big| R_{x_0} \Big( \frac{p_n}{q_n} + \epsilon h_n \Big)  - R_{x_0}( t ) \Big| + \frac{\sqrt{|h_n|}}{\sqrt{q_n}} \left(  1 + O\big( q_n^2 h_n \big) \right) + 2\pi (1 +  \epsilon) |h_n| \\
& \leq \Big| R_{x_0} \Big( \frac{p_n}{q_n} + \epsilon h_n \Big)  - R_{x_0}( t ) \Big| + 2\delta\, \frac{\sqrt{|h_n|}}{\sqrt{q_n}}.
\end{split}
\end{equation} 
In the last line we used the hypothesis of Case 1.2 and $|h_n| \leq \frac{\sqrt{|h_n|}}{\sqrt{q_n}} \, \frac{1}{\sqrt{q_n}}$.
Hence, 
 \begin{equation}
 \Big| R_{x_0} \Big( \frac{p_n}{q_n} + \epsilon h_n \Big)  - R_{x_0}( t ) \Big| 
 \gtrsim \left(  \frac{\sqrt{\epsilon}}{2}  - C\delta \right)   \frac{\sqrt{|h_n|}}{\sqrt{q_n}},
\end{equation}
for some $C>0$.
Fix $\sqrt{\epsilon} = 4C \delta$ small enough.
Writing $p_n/q_n + \epsilon h_n = t - (1-\epsilon) h_n$ and observing that $(1-\epsilon)|h_n| \simeq |h_n|$, we conclude that 
\begin{equation}
 \Big| R_{x_0} \big( t - (1- \epsilon) h_n \big)  - R_{x_0}( t ) \Big| 
 \gtrsim \delta \,  \frac{\sqrt{|h_n|}}{\sqrt{q_n}}
 \geq \delta \, |h_n|^{3/4} 
 \simeq | (1-\epsilon) h_n |^{3/4}, 
 \quad \text{ for large enough } n. 
\end{equation}
Hence $\alpha_{x_0}(t) \leq 3/4$. 
 
\end{itemize}
	
	
	\item[\textbf{Case 2}] We have $x_{q_n} \neq 0$ for all large enough $n \in \mathbb N$,
	hence $1/Q \leq q_n x_{q_n} \leq 1/2$. 
	We now use \eqref{eq:Asymptotic_xq_Not_0} which
	%, as opposed to \eqref{eq:Asymptotic_xq_Equal_0}, 
	has no leading $h^{1/2}$ term. 
	Rewrite it\footnote{
	When $q$ is odd and coprime with $p$, the inverses of 2 and $p$ modulo $q$ exist.
	Therefore, 
	\begin{equation}
	G(p,m,q) = \sum_{r=1}^q e^{2\pi i \frac{pr^2 + mr}{q}}
%	= \sum_{r=1}^q e^{2\pi i p \, \frac{r^2 + p^{-1}mr}{q}}
	= e^{2\pi i (4p)^{-1}\frac{m^2}{q}} \, \sum_{r=1}^q e^{2\pi i p \, \frac{(r + (2p)^{-1}m)^2}{q} } 
	= e^{2\pi i (4p)^{-1}\frac{m^2}{q}} \, G(p,0,q).
	\end{equation}	
	}, %\eqref{eq:Asymptotic_xq_Not_0} as
	assuming $1/Q \leq qx_q \leq 1/2$, 
	as	
	\begin{equation}
	\begin{split}
& R_{x_0}\Big( \frac{p}{q} + h \Big) - R_{x_0}\Big( \frac{p}{q} \Big) + 2\pi i h  \\
& \qquad \quad = 2(1 \pm i) \, \frac{G(p,0,q)}{\sqrt{q}}\, \frac{\sqrt{|h|}}{\sqrt{q}} \,  q^{2} |h| \, \left[  \sum_{m \in \mathbb Z} \, e^{2\pi i (4p)^{-1} \frac{(m_q + m)^2}{q}} \,  
\frac{e^{- 2\pi i \frac{(m - qx_q)^2}{4 q^2 h}}}{ (m - qx_q)^2} \,  + O_Q\left( q^{2} h  \right)
\right].
\end{split}
	\end{equation}
	Define the auxiliary function
	\begin{equation}\label{eq:f_q}
	f_q(y) = \sum_{m \in \mathbb Z} \, e^{2\pi i (4p)^{-1} \frac{ m^2 + 2m_q m}{q}} \,  
\frac{e^{- 2\pi i (m^2  - 2 qx_q m)  y }}{ (m - qx_q)^2}.
	\end{equation}
	Take absolute values and write
	\begin{equation}\label{eq:Asymptotic_f_q}
	\Big|  R_{x_0}\Big( \frac{p}{q} + h \Big) - R_{x_0}\Big( \frac{p}{q} \Big) + 2\pi i h  \Big|
	= 2\sqrt{2}  \frac{\sqrt{|h|}}{\sqrt{q}} \,  q^{2} |h| \,
\left|  f_q\Big( \frac{1}{4q^2 h}  \Big)  + O_Q\left( q^{2} h  \right)
\right|.
	\end{equation}
We now state the properties of this function, whose proof we postpone. 
\begin{lem}\label{thm:Lemma_f_q}
Let $q \in \mathbb N$, let $p \in \mathbb N$ be coprime with $q$ and $f_q$ defined in \eqref{eq:f_q}. 
Then, 
\begin{enumerate}
	\item $f_q$ is periodic of period $Q$.
	\item there exists $y_0^q \in [0,Q]$ depending on $q$ (and on $p$) such that $|f_q(y_0^q) | \geq 5$.
	\item The sequence defined by $y_k^q = y_0^q + kQ$ satisfies 
	\begin{equation}
	\lim_{k \to \infty} y_k^q = \infty, \qquad \text{ and } \qquad  |f_q(y_k^q)| \geq 5, \quad \forall k \in \mathbb N. 
	\end{equation}
\end{enumerate}
\end{lem}
\begin{rem}
The dependence on $p$ of the point $y_0^q$ is irrelevant for our purposes.
Indeed, once we fix $t \not\in \mathbb Q$, we get the sequence of approximations $p_n/q_n$, 
hence each $q_n$ comes with one and only one $p_n$. 
Hence, we can assume that the sequence $f_{q_n}$ only depends on $q_n$. 
\end{rem}
We now evaluate \eqref{eq:Asymptotic_f_q} at $p_n/q_n$ and $h_n = t - p_n/q_n$ and we separate two cases:
\begin{itemize}
	\item[\textbf{Case 2.1.}] Suppose $\limsup_{n \to \infty} q_n^2 |h_n| > 0$, so that there exists $c > 0$ and a subsequence for which $c < q_n^2 |h_n| \leq 1$.
	Then, from \eqref{eq:Asymptotic_f_q} we get
		\begin{equation}
	\Big|  R_{x_0}( t ) - R_{x_0}\Big( \frac{p_n}{q_n} \Big) + 2\pi i h_n  \Big|
	\geq c  \frac{\sqrt{|h_n|}}{\sqrt{q_n}} \,
\left|  f_{q_n}\Big( \frac{1}{4q_n^2 h_n}  \Big)  + O_Q\left( q_n^2 h_n  \right)
\right|.
	\end{equation}
	Fix $\delta>0$ which we later determine. Proceeding like in \textbf{Case 1}, we separate two cases:
	\begin{itemize}
		\item[\textbf{Case 2.1.1.}] Suppose $\left|  f_{q_n}\Big( \frac{1}{4q_n^2 h_n}  \Big)  + O_Q\left( q_n^2 h_n  \right) \right| \geq \delta$ for infinitely many $n$. Then, 
		\begin{equation}
	\Big|  R_{x_0}( t ) - R_{x_0}\Big( \frac{p_n}{q_n} \Big) + 2\pi i h_n  \Big|
	\geq c \delta \frac{\sqrt{|h_n|}}{\sqrt{q_n}} 
	\geq c\delta |h_n|^{3/4} 
	\end{equation}
		for infinitely many $n$, which implies $\alpha_{x_0}(t) \leq 3/4$.
				
		\item[\textbf{Case 2.1.2.}] Suppose $\left|  f_{q_n}\Big( \frac{1}{4q_n^2 h_n}  \Big)  + O_Q\left( q_n^2 h_n  \right) \right| < \delta$ for all large enough $n$. 
		Then, let $\epsilon_n$ be a sequence which we determine later,
		and define $\eta_n = \epsilon_n /q_n^2$. 
		Observe that $\eta_n = \epsilon_n |h_n| / (q_n^2 |h_n|) \simeq \epsilon_n |h_n|$. 
		Evaluate \eqref{eq:Asymptotic_f_q} at $\eta_n$ to get 
		\begin{equation}
		\begin{split}
		\Big|  R_{x_0}\Big( \frac{p_n}{q_n} + \eta_n \Big) - R_{x_0}\Big( \frac{p_n}{q_n} \Big) + 2\pi i \eta_n  \Big|
	& = 2\sqrt{2}  \frac{\sqrt{\eta_n}}{\sqrt{q_n}} \,  q_n^2 \eta_n \,
\left|  f_q\Big( \frac{1}{4q_n^2 \eta_n}  \Big)  + O_Q\left( q_n^{2} \eta_n  \right)
\right| \\
& = 2\sqrt{2} \,   \epsilon_n \,  \frac{\sqrt{\eta_n}}{\sqrt{q_n}} \,  
\left|  f_{q_n}\Big( \frac{1}{4 \epsilon_n}  \Big)  + O_Q\left( \epsilon_n  \right)
\right|. 
\end{split}
		\end{equation}
		Fix $k \in \mathbb N$ large enough and set $\epsilon_n = 1/(4y_k^{q_n})$.
		Then, by Lemma~\ref{thm:Lemma_f_q} (c), 
		\begin{equation}
		\Big|  f_{q_n}\Big( \frac{1}{4 \epsilon_n}  \Big) \Big| = \left|  f_{q_n}( y_K^{q_n}) \right| \geq 5,  \quad \forall n \text{ large enough.}
		\end{equation}
		Since $\epsilon_n \simeq 1/(kQ)$, 
		if $k \in \mathbb N$ is large enough we get $ O_Q\left( \epsilon_n  \right) \leq C_Q \epsilon_n \leq 1$.
		In particular, $|h_n| \simeq_Q k \eta_n$. 
		Therefore, 
		\begin{equation}
		\Big|  R_{x_0}\Big( \frac{p_n}{q_n} + \eta_n \Big) - R_{x_0}\Big( \frac{p_n}{q_n} \Big) + 2\pi i \eta_n  \Big|
 \geq  \epsilon_n\,  \frac{\sqrt{\eta_n}}{\sqrt{q_n}} 
 \simeq \epsilon_n^{3/2} \frac{\sqrt{|h_n|}}{\sqrt{q_n}}. 
		\end{equation}
		With this, and using the assumption of this case in \eqref{eq:Asymptotic_f_q}, we write
		\begin{equation}
		\begin{split}
		\epsilon_n^{3/2} \frac{\sqrt{|h_n|}}{\sqrt{q_n}}
		&  \lesssim \Big|  R_{x_0}\Big( \frac{p_n}{q_n} + \eta_n \Big) - R_{x_0}\Big( \frac{p_n}{q_n} \Big)\Big| + 2\pi \eta_n   \\
		& \leq \Big|  R_{x_0}\Big( \frac{p_n}{q_n} + \eta_n \Big) - R_{x_0}(t) \Big| 
		+ \Big| R_{x_0}(t) - R_{x_0}\Big( \frac{p_n}{q_n} \Big)\Big| 
		+ 2\pi \eta_n \\
		& \lesssim \Big|  R_{x_0}\Big( \frac{p_n}{q_n} + \eta_n \Big) - R_{x_0}(t) \Big| 
		+ \delta \frac{\sqrt{|h_n|}}{\sqrt{q_n}} 
		+ 2\pi ( |h_n|  + \eta_n)   \\
		& \lesssim \Big|  R_{x_0}\Big( \frac{p_n}{q_n} + \eta_n \Big) - R_{x_0}(t) \Big| 
		+ \delta \frac{\sqrt{|h_n|}}{\sqrt{q_n}},
		\end{split}
		\end{equation}
	for large enough $n$, 
	where in the last line we used $\eta_n \simeq |h_n|/k \leq |h_n|$
	and $|h_n| \leq  \frac{\sqrt{|h_n|}}{\sqrt{q_n}} \, \frac{1}{\sqrt{q_n}} \ll  \frac{\sqrt{|h_n|}}{\sqrt{q_n}}$. 
	Since $\epsilon_n \simeq_Q 1/k$, set  $\delta = 1/(c_Qk^{3/2})$ with some small enough $c_Q>0$ so that
	\begin{equation}
	\Big|  R_{x_0}\Big( \frac{p_n}{q_n} + \eta_n \Big) - R_{x_0}(t) \Big| 
	\gtrsim \Big(  \epsilon_n^{3/2} - C\delta \Big) \frac{\sqrt{|h_n|}}{\sqrt{q_n}} 
	\gtrsim \delta \frac{\sqrt{|h_n|}}{\sqrt{q_n}} 
	\geq \delta  |h_n|^{3/4}. 
	\end{equation}
	Write $p_n/q_n + \eta_n = t - (h_n -  \eta_n)$. 
	Since $|h_n - \eta_n| \leq 2|h_n|$, 
	we get 
	\begin{equation}
	\Big|  R_{x_0}\Big(  t - (h_n -  \eta_n) \Big) - R_{x_0}(t) \Big| 
	\geq \delta |h_n|^{3/4} 
	\gtrsim \delta |h_n -  \eta_n|^{3/4}, \quad \text{ for large enough } n, 
	\end{equation}
	which implies $\alpha_{x_0}(t) \leq 3/4$. 
		
	\end{itemize}
	
	
	\item[\textbf{Case 2.2.}] Suppose $\lim_{n \to \infty} q_n^2 |h_n| = 0$. 
	In this case, the term $q_n^2 |h_n|$ in \eqref{eq:Asymptotic_f_q} tends to zero, which kills the desired $|h_n|^{3/4}$ that came from $\sqrt{h_n}/\sqrt{q_n}$. 
	To counter that, define $\eta_n = \epsilon_n / q_n^2$ 
	as in Case 2.1.2. By \eqref{eq:Asymptotic_f_q}, 
	\begin{equation}
	\Big|  R_{x_0}\Big( \frac{p_n}{q_n} + \eta_n \Big) - R_{x_0}\Big( \frac{p_n}{q_n} \Big) + 2\pi i \eta_n  \Big|
	= 2\sqrt{2} \epsilon_n \frac{\sqrt{\eta_n}}{\sqrt{q_n}} \,
\left|  f_{q_n}\Big( \frac{1}{4\epsilon_n}  \Big)  + O_Q\left(\epsilon_n  \right)
\right|.
	\end{equation}
	Fix $k \in \mathbb N$ large enough and set $\epsilon_n = 1/(4y_k^{q_n})$. Then,
	\begin{equation}
	\left|  f_{q_n}\Big( \frac{1}{4\epsilon_n}  \Big)  \right| 
	= \left|  f_{q_n}\Big( y_k^{q_n} \Big)  \right| \geq 5, 
	\quad \text{ and } \quad
	O_Q(\epsilon_n) \leq C_Q \epsilon_n = \frac{C_Q}{4y_k^{q_n}} \simeq \frac{C_Q}{kQ} \leq 1,
	\end{equation}
	so
	\begin{equation}
	\Big|  R_{x_0}\Big( \frac{p_n}{q_n} + \eta_n \Big) - R_{x_0}\Big( \frac{p_n}{q_n} \Big) + 2\pi i \eta_n  \Big|
	\geq  \epsilon_n \frac{\sqrt{\eta_n}}{\sqrt{q_n}}. 
	\end{equation}
	With this and \eqref{eq:Asymptotic_f_q}, we can write
	\begin{equation}
	\begin{split}
	\epsilon_n \frac{\sqrt{\eta_n}}{\sqrt{q_n}} 
	&  \leq \Big|  R_{x_0}\Big( \frac{p_n}{q_n} + \eta_n \Big) - R_{x_0}\Big( \frac{p_n}{q_n} \Big) \Big| + 2\pi \eta_n  \\
	 &  \leq \Big|  R_{x_0}\Big( \frac{p_n}{q_n} + \eta_n \Big) -R_{x_0}(t) \Big| + \Big| R_{x_0}(t) -  R_{x_0}\Big( \frac{p_n}{q_n} \Big) \Big| + 2\pi \eta_n  \\ 
	 &  \lesssim \Big|  R_{x_0}\Big( \frac{p_n}{q_n} + \eta_n \Big) -R_{x_0}(t) \Big| + \frac{\sqrt{|h_n|}}{\sqrt{q_n}} \, q_n^2 |h_n| + 2\pi (\eta_n + |h_n|).
	 \end{split}
	\end{equation}
	Since $\lim_{n \to \infty} q_n^2 |h_n|=0$ implies $h_n = o(\eta_n)$, and $\eta_n =  \frac{\sqrt{\eta_n}}{\sqrt{q_n}} \frac{\sqrt{\epsilon_n}}{\sqrt{q_n}}$, 
	we get
	\begin{equation}
	\Big|  R_{x_0}\Big( \frac{p_n}{q_n} + \eta_n \Big) -R_{x_0}(t) \Big|
	\gtrsim \Big( \epsilon_n - q_n^2 h_n - \frac{\sqrt{\epsilon_n}}{ \sqrt{q_n} } \Big) \frac{\sqrt{\eta_n}}{\sqrt{q_n}}
	\geq \frac{\epsilon_n}{2} \, \frac{\sqrt{\eta_n}}{\sqrt{q_n}}
	=  \frac{\epsilon_n^{3/4}}{2} \, \eta_n^{3/4}.
	\end{equation}
	Write $p_n/q_n + \eta_n = t + (\eta_n - h_n)$. 
	Recalling $\epsilon_n \simeq 1/(kQ)$ for all $n$, and since $h_n = o(\eta_n)$ implies $|\eta_n - h_n| \simeq \eta_n$,  we conclude
	\begin{equation}
	\Big|  R_{x_0}\Big( t + (\eta_n - h_n) \Big) -R_{x_0}(t) \Big|
	\geq  \frac{\epsilon_n^{3/4}}{2} \, \eta_n^{3/4}
	\simeq_Q |\eta_n - h_n|^{3/4}, 
	\end{equation}
	and therefore $\alpha_{x_0}(t) \leq 3/4$. 
\end{itemize}

\end{itemize}
\end{proof}

\end{document}
